% Generated draft: requires mathematical and source review.
\documentclass[11pt]{article}
\usepackage[margin=25mm]{geometry}
\usepackage{fontspec}
\setmainfont{Latin Modern Roman}
\usepackage{amsmath,amssymb,mathtools,unicode-math}
\setmathfont{Latin Modern Math}
\usepackage{enumitem,needspace,xurl,hyperref}
\setlength{\parindent}{0pt}
\setlength{\parskip}{5pt}
\setlength{\emergencystretch}{3em}
\newcommand{\N}{\mathbb{N}}
\newcommand{\Z}{\mathbb{Z}}
\newcommand{\Q}{\mathbb{Q}}
\newcommand{\R}{\mathbb{R}}
\newcommand{\C}{\mathbb{C}}
\newcommand{\F}{\mathbb{F}}
\newcommand{\A}{\mathbb{A}}
\newcommand{\PP}{\mathbb{P}}
\newcommand{\E}{\mathbb{E}}
\newcommand{\eps}{\varepsilon}
\newcommand{\dd}{\,\mathrm{d}}
\begin{document}
\section*{UnsolvedMath ID: 20002828}
\textbf{Source problem code:} catalogue entry
\subsection{Definitions and mathematical statement}
The source asks whether the Tung Le--Kay Magaard method of refining supercharacters of \(U_n\) has an analogue for arbitrary algebra groups. Its remark points to T. Le's paper \emph{Supercharacters and pattern subgroups in the upper triangular groups}, without further bibliographic details. The exact refinement method and desired terminal output are not described in the source.

Fix a finite-dimensional nilpotent associative \(\F_q\)-algebra \(J\), the group \(G=1+J\), a nontrivial additive character \(\psi:\F_q\to\C^*\), and \(\lambda\in J^*\). Let \(\chi_\lambda\) be the standard Diaconis--Isaacs supercharacter, normalized as the induction of the linear character below. Define
\[
 \mathfrak l_\lambda=\{x\in J:\lambda(xy)=0\text{ for all }y\in J\},
\]
\[
 \mathfrak s_\lambda=\{x\in J:\lambda(xy)=0\text{ for all }y\in\mathfrak l_\lambda\},
 \qquad L_\lambda=1+\mathfrak l_\lambda,\quad
 S_\lambda=1+\mathfrak s_\lambda,
\]
\[
 \vartheta_\lambda(1+x)=\psi(\lambda(x))\quad(x\in\mathfrak l_\lambda),
 \qquad\chi_\lambda=\operatorname{Ind}_{L_\lambda}^G\vartheta_\lambda.
\]
These definitions fix the left-kernel convention; reversing left/right conventions requires corresponding changes throughout.

The aid reports the 2008 Andr\'e--Nicolas constituent correspondence and Marberg's iterative development as the natural algebra-group analogue: supercharacters can be refined through canonical algebra subgroups, with induction preserving irreducibility and distinguishing the relevant induced constituents. The general result is a constituent correspondence/refinement mechanism, not a supplied closed classification of all algebra-group irreducibles or a guarantee that no residual subgroup problem remains. The exact general theorem and iterative stopping hypotheses are not reproduced in the record.

\textbf{Closed class-two formula reported in the aid.} Suppose \(J^3=0\). Then \(\mathfrak l_\lambda\) is an ideal in \(\mathfrak s_\lambda\), \(S_\lambda/L_\lambda\) is the additive vector space \(\mathfrak s_\lambda/\mathfrak l_\lambda\), and the alternating form
\[
 \beta_\lambda(x+\mathfrak l_\lambda,y+\mathfrak l_\lambda)
             =\lambda(xy-yx)
\]
is well defined. Let its radical have \(\F_q\)-dimension \(r\) and its rank be \(2m\), so \(\dim(\mathfrak s_\lambda/\mathfrak l_\lambda)=r+2m\). The form is alternating also in characteristic two; no odd-characteristic assumption is inserted.

The aid reports that the irreducible constituents of \(\chi_\lambda\) are exactly \(q^r\) distinct characters, each of degree \([G:S_\lambda]q^m\), each occurring with multiplicity \(q^m\). Equivalently,
\[
 \chi_\lambda=q^m\sum_{a=1}^{q^r}\rho_a,\qquad
 \rho_a\ne\rho_b\ (a\ne b),\qquad
 \rho_a(1)=[G:S_\lambda]q^m.
\]
The degree check is
\[
 q^r q^m[G:S_\lambda]q^m
        =[G:S_\lambda]q^{r+2m}=[G:L_\lambda].
\]
This gives number, common degree and common multiplicity of constituents, not a canonical labeling of all \(\rho_a\). It is a special-case explicit decomposition for \(J^3=0\), while the general refinement question is addressed by the reported subgroup correspondence. The supplied solved-in-literature label is not an independent theorem verification or novelty claim.
\subsection{Short English statement}
The aid reports an algebra-group refinement correspondence extending the unitriangular idea through canonical subgroups. When the algebra has cube zero, an alternating form gives an explicit decomposition: its radical and rank determine the number, degree and multiplicity of all constituents. This special formula does not amount to a closed classification for arbitrary nilpotent algebras.
\subsection{Sources}
\begin{itemize}
\item Local UnsolvedMath registry, permanent ID 20002828, catalogue entry; canonical \texttt{aim-representation-theory-notes.json}, notes[99].
\item AIM \emph{Supercharacters and combinatorial Hopf algebras}, Problem 1.16, \url{https://aimath.org/WWN/supercharacters/supercharacters.pdf}. Source names Tung Le and Kay Magaard.
\item Original remark: T. Le, \emph{Supercharacters and pattern subgroups in the upper triangular groups}; no year, journal or URL is supplied.
\item Aid attributions: 2008 Andr\'e--Nicolas constituent correspondence; Marberg's iterative refinement; Diaconis--Isaacs supercharacters. Full references and exact general theorem hypotheses are absent.
\item Machine-generated research attempt 1, partial-result classification, reported problem status solved, rigor proved and low novelty confidence; registry classification solved in literature. The class-two formula and refinement reports are not peer reviewed and require independent verification.
\end{itemize}
\end{document}
